\documentclass{ximera}
%verbeterd 9/9/26
\title{Theorie: Elektrische krachtwerking}
\author{Daan Lipkens}

\begin{document}
\begin{abstract}
    Inleiding tot ladingen en de elektrische krachtwerking daartussen.
\end{abstract}

\maketitle
\label{hoofdstuk:elektrische_krachtwerking}
% ----------------------------------------------------------------------
%          inleiding hoofdstuk
% ----------------------------------------------------------------------
\section*{Inleiding}
De elektromagnetische kracht tussen geladen deeltjes is een van de vier fundamentele natuurkrachten. In dit hoofdstuk onderzoeken we eerst de basiseigenschappen van elektrische lading en de krachten die erop inwerken.

Daarna verdiepen we ons in de \emph{wet van Coulomb}, die zowel de grootte als de richting van de kracht tussen geladen deeltjes beschrijft. Je leert hoe je deze kracht kunt berekenen in uiteenlopende situaties.

Dit geheel behoort tot de \emph{elektrostatica}: het gedrag van stilstaande ladingen. In latere hoofdstukken bestuderen we ook bewegende ladingen en de verschijnselen die daarbij horen.

\begin{definition}
\textbf{Elektrostatica} is het deel van de fysica dat het gedrag en de interacties van elektrische ladingen in rust en de krachtwerking ertussen bestudeert.
\end{definition}

\end{document}